\documentclass[a4paper,12pt]{article}

\usepackage{alltt, fancyvrb, url}
\usepackage{graphicx}
\usepackage[utf8]{inputenc}
\usepackage{float}
\usepackage{hyperref}
\usepackage[italian]{babel}
\usepackage[table]{xcolor}
\usepackage{array}
\usepackage{multirow}
\usepackage{titlesec}
\usepackage{listings}
\usepackage[italian]{cleveref}
\usepackage{acronym}

\renewcommand{\thesection}{\arabic{section}}
\renewcommand{\labelenumii}{\arabic{enumi}.\arabic{enumii}}
\renewcommand{\labelenumiii}{\arabic{enumi}.\arabic{enumii}.\arabic{enumiii}}
\renewcommand{\labelenumiv}{\arabic{enumi}.\arabic{enumii}.\arabic{enumiii}.\arabic{enumiv}}

\sloppy
\setlength{\parindent}{0pt}

\title{\textbf{Prioritizzazione dei requisiti}}
\author{}
\date{}

\begin{document}
\maketitle

In riferimento al RBS definito, prioritizziamo i requisiti in due categorie: \textit{Must} e \textit{Nice to have}. I
requisiti \textit{Must} sono quelli che devono essere implementati per garantire il successo del progetto, mentre i
requisiti \textit{Nice to have} sono quelli che possono essere implementati se le risorse e il tempo lo permettono, ma
non sono essenziali per il successo del progetto. Nella \cref{table:requirements_prioritized} sono elencati i requisiti
delle foglie del RBS secondo le due categorie.

\begin{table}[!htb]
    \centering
    \begin{tabular}{|c|c|}
        \hline
        Must  & Nice to have \\
        \hline
        1.1.1 & 1.4          \\
        1.1.2 &              \\
        1.1.3 &              \\
        1.2   &              \\
        1.3   &              \\
        1.5   &              \\
        2.1   &              \\
        2.2   &              \\
        2.3   &              \\
        3.1   &              \\
        3.2   &              \\
        3.3   &              \\
        4.1   &              \\
        4.2   &              \\
        5.1   &              \\
        5.2   &              \\    
        \hline
    \end{tabular}
    \caption{Prioritizzazione dei requisiti del RBS.}
    \label{table:requirements_prioritized}
\end{table}

\end{document}